\section{Details: Limits of FPSs}

\subsection{Detailed proof of Lemma~\ref{lem.fps.lim.xn-equiv}}

\begin{proof}[Detailed proof of Lemma~\ref{lem.fps.lim.xn-equiv}]
\proves{lem.fps.lim.xn-equiv}
\leanok
We have $\lim_{i\to\infty} f_i = f$.
In other words, the sequence $(f_i)_{i\in\mathbb{N}}$ coefficientwise
stabilizes to $f$.
In other words, for each $n\in\mathbb{N}$,
\begin{equation}
\text{the sequence } \bigl([x^n] f_i\bigr)_{i\in\mathbb{N}}
\text{ stabilizes to } [x^n] f
\label{pf.lem.fps.lim.xn-equiv.stab}
\end{equation}
(by the definition of ``coefficientwise stabilizing'').

Now, let $n\in\mathbb{N}$. Furthermore, let $k\in\{0,1,\ldots,n\}$.
Then, the sequence $\bigl([x^k] f_i\bigr)_{i\in\mathbb{N}}$ stabilizes
to $[x^k] f$ (by~\eqref{pf.lem.fps.lim.xn-equiv.stab}, applied to $k$
instead of $n$). In other words, there exists some $N\in\mathbb{N}$
such that
\[
\text{all integers } i\geq N \text{ satisfy } [x^k] f_i = [x^k] f
\]
(by the definition of ``stabilizing''). Let us denote this $N$ by $N_k$.
Thus,
\begin{equation}
\text{all integers } i\geq N_k \text{ satisfy } [x^k] f_i = [x^k] f.
\label{pf.lem.fps.lim.xn-equiv.Nk}
\end{equation}

Forget that we fixed $k$. Thus, for each $k\in\{0,1,\ldots,n\}$, we
have defined an integer $N_k\in\mathbb{N}$ for
which~\eqref{pf.lem.fps.lim.xn-equiv.Nk} holds. Altogether, we have
thus defined $n+1$ integers $N_0, N_1, \ldots, N_n \in \mathbb{N}$.

Let us set $P := \max\{N_0, N_1, \ldots, N_n\}$. Then, of course,
$P\in\mathbb{N}$.

Now, let $i\geq P$ be an integer. Then, for each $k\in\{0,1,\ldots,n\}$,
we have $i\geq P = \max\{N_0, N_1, \ldots, N_n\} \geq N_k$ (since
$k\in\{0,1,\ldots,n\}$) and therefore $[x^k] f_i = [x^k] f$
(by~\eqref{pf.lem.fps.lim.xn-equiv.Nk}). In other words, each
$k\in\{0,1,\ldots,n\}$ satisfies $[x^k] f_i = [x^k] f$. Renaming the
variable $k$ as $m$ in this result, we obtain the following:
\[
\text{Each } m\in\{0,1,\ldots,n\} \text{ satisfies } [x^m] f_i = [x^m] f.
\]
In other words, $f_i \overset{x^n}{\equiv} f$ (by the definition of
$x^n$-equivalence).

Forget that we fixed $i$. We thus have shown that all integers $i\geq P$
satisfy $f_i \overset{x^n}{\equiv} f$. Hence, there exists some integer
$N\in\mathbb{N}$ such that
\[
\text{all integers } i\geq N \text{ satisfy }
f_i \overset{x^n}{\equiv} f
\]
(namely, $N = P$). This proves Lemma~\ref{lem.fps.lim.xn-equiv}.
\end{proof}

\subsection{Detailed proof of Proposition~\ref{prop.fps.lim.sum-prod}}

\begin{lemma}
\label{prop.fps.lim.sum-prod.K}
\lean{PowerSeries.exists_xnEquiv_K}
\leantarget
\leanok
Let $(f_i)_{i\in\mathbb{N}}$ be a sequence of FPSs that coefficientwise
stabilizes to an FPS $f$, and let $n\in\mathbb{N}$. Then there exists
an integer $K\in\mathbb{N}$ such that
\[
\text{all integers } i\geq K \text{ satisfy }
f_i \overset{x^n}{\equiv} f.
\]
\end{lemma}

\begin{proof}
\leanok
This follows directly from Lemma~\ref{lem.fps.lim.xn-equiv}
applied to the sequence $(f_i)$ and limit $f$.
\end{proof}

\begin{lemma}
\label{prop.fps.lim.sum-prod.L}
\lean{PowerSeries.exists_xnEquiv_L}
\leantarget
\leanok
Let $(g_i)_{i\in\mathbb{N}}$ be a sequence of FPSs that coefficientwise
stabilizes to an FPS $g$, and let $n\in\mathbb{N}$. Then there exists
an integer $L\in\mathbb{N}$ such that
\[
\text{all integers } i\geq L \text{ satisfy }
g_i \overset{x^n}{\equiv} g.
\]
\end{lemma}

\begin{proof}
\leanok
This is the exact same statement as
Lemma~\ref{prop.fps.lim.sum-prod.K}, but applied to the sequence
$(g_i)$ and limit $g$. It follows directly from
Lemma~\ref{lem.fps.lim.xn-equiv}.
\end{proof}

\begin{lemma}
\label{prop.fps.lim.sum-prod.both}
\lean{PowerSeries.exists_xnEquiv_both}
\leanhelper
\leanok
Let $(f_i)_{i\in\mathbb{N}}$ and $(g_i)_{i\in\mathbb{N}}$ be sequences
of FPSs that coefficientwise stabilize to $f$ and $g$ respectively, and
let $n\in\mathbb{N}$. Then there exists an integer $P\in\mathbb{N}$
such that for all $i\geq P$,
\[
f_i \overset{x^n}{\equiv} f
\qquad\text{and}\qquad
g_i \overset{x^n}{\equiv} g.
\]
\end{lemma}

\begin{proof}
\leanok
By Lemma~\ref{prop.fps.lim.sum-prod.K}, there exists $K$ such that
$f_i \overset{x^n}{\equiv} f$ for all $i\geq K$.
By Lemma~\ref{prop.fps.lim.sum-prod.L}, there exists $L$ such that
$g_i \overset{x^n}{\equiv} g$ for all $i\geq L$.
Set $P := \max\{K, L\}$. Then for all $i\geq P$, both
$i\geq K$ and $i\geq L$ hold, so both equivalences hold.
\end{proof}

\begin{proof}[Detailed proof of Proposition~\ref{prop.fps.lim.sum-prod}]
\proves{prop.fps.lim.sum-prod}
\leanok
Recall that $\lim_{i\to\infty} f_i = f$.
Hence, Lemma~\ref{lem.fps.lim.xn-equiv} shows that there exists some
integer $N\in\mathbb{N}$ such that
\[
\text{all integers } i\geq N \text{ satisfy }
f_i \overset{x^n}{\equiv} f.
\]
Let us denote this $N$ by $K$. Hence, all integers $i\geq K$ satisfy
$f_i \overset{x^n}{\equiv} f$.

Thus, we have found an integer $K\in\mathbb{N}$ such that
\[
\text{all integers } i\geq K \text{ satisfy }
f_i \overset{x^n}{\equiv} f.
\]
Similarly, using $\lim_{i\to\infty} g_i = g$, we can find an integer
$L\in\mathbb{N}$ such that
\[
\text{all integers } i\geq L \text{ satisfy }
g_i \overset{x^n}{\equiv} g.
\]
Consider these $K$ and $L$. Let us furthermore set
$P := \max\{K, L\}$. Thus, $P\in\mathbb{N}$.

We shall now show that each integer $i\geq P$ satisfies
$[x^n](f_i g_i) = [x^n](fg)$.

Indeed, let $i\geq P$ be an integer. Then, $f_i \overset{x^n}{\equiv} f$
(since $i\geq P = \max\{K, L\} \geq K$), whereas
$g_i \overset{x^n}{\equiv} g$
(since $i\geq P = \max\{K, L\} \geq L$). Hence, we obtain
\[
f_i g_i \overset{x^n}{\equiv} fg
\]
(by the multiplicativity of $x^n$-equivalence, i.e., Theorem~\ref{thm.fps.xneq.props}~\textbf{(b)}). In other words,
\[
\text{each } m\in\{0,1,\ldots,n\} \text{ satisfies }
[x^m](f_i g_i) = [x^m](fg)
\]
(by the definition of $x^n$-equivalence). Applying this to $m = n$,
we find
\[
[x^n](f_i g_i) = [x^n](fg).
\]

Forget that we fixed $i$. We thus have shown that all integers $i\geq P$
satisfy $[x^n](f_i g_i) = [x^n](fg)$. Hence, there exists some
$N\in\mathbb{N}$ such that
\[
\text{all integers } i\geq N \text{ satisfy }
[x^n](f_i g_i) = [x^n](fg)
\]
(namely, $N = P$). In other words,
\[
\text{the sequence } \bigl([x^n](f_i g_i)\bigr)_{i\in\mathbb{N}}
\text{ stabilizes to } [x^n](fg)
\]
(by the definition of ``stabilizes'').

Forget that we fixed $n$. We thus have shown that for each
$n\in\mathbb{N}$, the sequence $\bigl([x^n](f_i g_i)\bigr)_{i\in\mathbb{N}}$
stabilizes to $[x^n](fg)$. In other words, the sequence
$(f_i g_i)_{i\in\mathbb{N}}$ coefficientwise stabilizes to $fg$
(by the definition of ``coefficientwise stabilizing''). In other words,
$\lim_{i\to\infty}(f_i g_i) = fg$.

The same argument---but using the additivity of $x^n$-equivalence instead
of the multiplicativity---shows that
$\lim_{i\to\infty}(f_i + g_i) = f + g$. Thus, the proof of
Proposition~\ref{prop.fps.lim.sum-prod} is complete.
\end{proof}

\begin{lemma}
\label{prop.fps.lim.sum-prod.add-detail}
\lean{PowerSeries.coeffStabilizesTo_add'}
\leanhelper
\leanok
If $(f_i)$ coefficientwise stabilizes to $f$ and $(g_i)$ coefficientwise
stabilizes to $g$, then $(f_i + g_i)$ coefficientwise stabilizes to
$f + g$.
\end{lemma}

\begin{proof}
\leanok
For each $n\in\mathbb{N}$:
\begin{enumerate}
\item By Lemma~\ref{lem.fps.lim.xn-equiv}, obtain $K_f$ such that
  $f_i \overset{x^n}{\equiv} f$ for $i\geq K_f$, and $K_g$ such that
  $g_i \overset{x^n}{\equiv} g$ for $i\geq K_g$.
\item Set $P := \max\{K_f, K_g\}$.
\item For $i\geq P$, by the additivity of $x^n$-equivalence,
  $f_i + g_i \overset{x^n}{\equiv} f + g$.
\item Extracting the $n$-th coefficient:
  $[x^n](f_i + g_i) = [x^n](f + g)$.
\end{enumerate}
\end{proof}

\begin{lemma}
\label{prop.fps.lim.sum-prod.mul-detail}
\lean{PowerSeries.coeffStabilizesTo_mul'}
\leanhelper
\leanok
If $(f_i)$ coefficientwise stabilizes to $f$ and $(g_i)$ coefficientwise
stabilizes to $g$, then $(f_i \cdot g_i)$ coefficientwise stabilizes to
$f \cdot g$.
\end{lemma}

\begin{proof}
\leanok
For each $n\in\mathbb{N}$:
\begin{enumerate}
\item By Lemma~\ref{lem.fps.lim.xn-equiv}, obtain $K_f$ such that
  $f_i \overset{x^n}{\equiv} f$ for $i\geq K_f$, and $K_g$ such that
  $g_i \overset{x^n}{\equiv} g$ for $i\geq K_g$.
\item Set $P := \max\{K_f, K_g\}$.
\item For $i\geq P$, by the multiplicativity of $x^n$-equivalence,
  $f_i \cdot g_i \overset{x^n}{\equiv} f \cdot g$.
\item Extracting the $n$-th coefficient:
  $[x^n](f_i \cdot g_i) = [x^n](f \cdot g)$.
\end{enumerate}
\end{proof}

\subsection{Helpers for Proposition~\ref{prop.fps.lim.sum-quot}}

\begin{lemma}
\label{lem.fps.lim.xnEquiv-invOfUnit-prime}
\lean{PowerSeries.xnEquiv_invOfUnit'}
\leanhelper
\leanok
Let $f, g \in K[[x]]$ with $f \overset{x^n}{\equiv} g$, and let
$u, v$ be units in $K$ with $u = v$ as elements. Then
the formal inverse of $f$ with respect to $u$ is $x^n$-equivalent to
the formal inverse of $g$ with respect to $v$.
\end{lemma}

\begin{lemma}
\label{lem.fps.lim.xnEquiv-div}
\lean{PowerSeries.xnEquiv_div}
\leanhelper
\leanok
Let $a \overset{x^n}{\equiv} b$ and $c \overset{x^n}{\equiv} d$,
where $c$ and $d$ are invertible with constant terms $u$ and $v$
respectively. Then
$a \cdot c^{-1} \overset{x^n}{\equiv}
b \cdot d^{-1}$.
\end{lemma}

\begin{lemma}
\label{lem.fps.lim.unit-eq-of-coeffStabilizesTo}
\lean{PowerSeries.unit_eq_of_coeffStabilizesTo}
\leanhelper
\leanok
If $(g_i)$ coefficientwise stabilizes to $g$ and each $g_i$ has
invertible constant term, then eventually the unit associated to the
constant term of $g_i$ equals the unit associated to the constant term
of $g$.
\end{lemma}

\begin{lemma}
\label{lem.fps.lim.coeffStabilizesTo-invOfUnit}
\lean{PowerSeries.coeffStabilizesTo_invOfUnit}
\leanhelper
\leanok
If $(g_i)$ coefficientwise stabilizes to $g$ and each $g_i$ has
invertible constant term, then the sequence of formal inverses
$(g_i^{-1})_{i \in \mathbb{N}}$
coefficientwise stabilizes to $g^{-1}$. The proof
proceeds by strong induction on the coefficient index, using the
recursive formula for coefficients of the formal inverse.
\end{lemma}

\subsection{Detailed proof of Proposition~\ref{prop.fps.lim.sum-quot}}

\begin{proof}[Detailed proof of Proposition~\ref{prop.fps.lim.sum-quot}]
\proves{prop.fps.lim.sum-quot}
\leanok
First, let us show that $g$ is invertible:

\textbf{Claim 1:} \textit{The FPS $g\in K[[x]]$ is invertible.}

\begin{proof}[Proof of Claim 1]
\leanok
We have assumed that $\lim_{i\to\infty} g_i = g$. In other words, the
sequence $(g_i)_{i\in\mathbb{N}}$ coefficientwise stabilizes to $g$.
In other words, for each $n\in\mathbb{N}$, the sequence
$\bigl([x^n] g_i\bigr)_{i\in\mathbb{N}}$ stabilizes to $[x^n] g$
(by the definition of ``coefficientwise stabilizing''). Applying this to
$n = 0$, we see that the sequence $\bigl([x^0] g_i\bigr)_{i\in\mathbb{N}}$
stabilizes to $[x^0] g$. In other words, there exists some
$N\in\mathbb{N}$ such that
\[
\text{all integers } i\geq N \text{ satisfy }
[x^0] g_i = [x^0] g
\]
(by the definition of ``stabilizes''). Consider this $N$. Thus, all
integers $i\geq N$ satisfy $[x^0] g_i = [x^0] g$. Applying this to
$i = N$, we obtain $[x^0] g_N = [x^0] g$.

However, each FPS $g_i$ is invertible (by assumption). Hence, in
particular, the FPS $g_N$ is invertible. By
Proposition~\ref{prop.fps.invertible} (applied to $a = g_N$), this
entails that its constant term $[x^0] g_N$ is invertible in $K$. In
other words, $[x^0] g$ is invertible in $K$ (since
$[x^0] g_N = [x^0] g$).

But the FPS $g$ is invertible if and only if its constant term $[x^0] g$
is invertible in $K$ (by Proposition~\ref{prop.fps.invertible}, applied
to $a = g$). Hence, $g$ is invertible (since its constant term $[x^0] g$
is invertible in $K$). This proves Claim 1.
\end{proof}

It remains to prove that
$\lim_{i\to\infty} \frac{f_i}{g_i} = \frac{f}{g}$.

Recall that $\lim_{i\to\infty} f_i = f$. Hence,
Lemma~\ref{lem.fps.lim.xn-equiv} shows that there exists some integer
$N\in\mathbb{N}$ such that
\[
\text{all integers } i\geq N \text{ satisfy }
f_i \overset{x^n}{\equiv} f.
\]
Let us denote this $N$ by $K$. Hence, all integers $i\geq K$ satisfy
$f_i \overset{x^n}{\equiv} f$.

Thus, we have found an integer $K\in\mathbb{N}$ such that
\[
\text{all integers } i\geq K \text{ satisfy }
f_i \overset{x^n}{\equiv} f.
\]
Similarly, using $\lim_{i\to\infty} g_i = g$, we can find an integer
$L\in\mathbb{N}$ such that
\[
\text{all integers } i\geq L \text{ satisfy }
g_i \overset{x^n}{\equiv} g.
\]
Consider these $K$ and $L$. Let us furthermore set
$P := \max\{K, L\}$. Thus, $P\in\mathbb{N}$.

We shall now show that each integer $i\geq P$ satisfies
$[x^n] \frac{f_i}{g_i} = [x^n] \frac{f}{g}$.

Indeed, let $i\geq P$ be an integer. Then,
$f_i \overset{x^n}{\equiv} f$
(since $i\geq P = \max\{K, L\} \geq K$), whereas
$g_i \overset{x^n}{\equiv} g$
(since $i\geq P = \max\{K, L\} \geq L$). Since both FPSs $g_i$ and
$g$ are invertible (by Claim 1), we thus obtain
\[
\frac{f_i}{g_i} \overset{x^n}{\equiv} \frac{f}{g}
\]
(by Theorem~\ref{thm.fps.xneq.props}~\textbf{(e)}, applied to
$a = f_i$, $b = f$, $c = g_i$ and $d = g$). In other words,
\[
\text{each } m\in\{0,1,\ldots,n\} \text{ satisfies }
[x^m] \frac{f_i}{g_i} = [x^m] \frac{f}{g}
\]
(by the definition of $x^n$-equivalence). Applying this to $m = n$,
we find
\[
[x^n] \frac{f_i}{g_i} = [x^n] \frac{f}{g}.
\]

Forget that we fixed $i$. We thus have shown that all integers $i\geq P$
satisfy $[x^n] \frac{f_i}{g_i} = [x^n] \frac{f}{g}$. Hence, there
exists some $N\in\mathbb{N}$ such that
\[
\text{all integers } i\geq N \text{ satisfy }
[x^n] \frac{f_i}{g_i} = [x^n] \frac{f}{g}
\]
(namely, $N = P$). In other words,
\[
\text{the sequence } \left([x^n] \frac{f_i}{g_i}\right)_{i\in\mathbb{N}}
\text{ stabilizes to } [x^n] \frac{f}{g}
\]
(by the definition of ``stabilizes'').

Forget that we fixed $n$. We thus have shown that for each
$n\in\mathbb{N}$, the sequence
$\left([x^n] \frac{f_i}{g_i}\right)_{i\in\mathbb{N}}$ stabilizes to
$[x^n] \frac{f}{g}$. In other words, the sequence
$\left(\frac{f_i}{g_i}\right)_{i\in\mathbb{N}}$ coefficientwise
stabilizes to $\frac{f}{g}$ (by the definition of ``coefficientwise
stabilizing''). In other words,
$\lim_{i\to\infty} \frac{f_i}{g_i} = \frac{f}{g}$.
This proves Proposition~\ref{prop.fps.lim.sum-quot} (since we already
have shown that $g$ is invertible).
\end{proof}

\subsection{Detailed proof of Theorem~\ref{thm.fps.lim.sum-lim}}

\begin{proof}[Proof of Theorem~\ref{thm.fps.lim.sum-lim}]
\proves{thm.fps.lim.sum-lim}
\leanok
The family $(f_n)_{n\in\mathbb{N}}$ is summable. In other words,
the family $(f_k)_{k\in\mathbb{N}}$ is summable (since this is the same
family as $(f_n)_{n\in\mathbb{N}}$). In other words, for each
$n\in\mathbb{N}$,
\begin{equation}
\text{all but finitely many } k\in\mathbb{N} \text{ satisfy }
[x^n] f_k = 0
\label{pf.thm.fps.lim.sum-lim.1}
\end{equation}
(by the definition of ``summable'').

Let us define
\begin{equation}
g_i := \sum_{k=0}^{i} f_k
\qquad\text{for each } i\in\mathbb{N}.
\label{pf.thm.fps.lim.sum-lim.gi}
\end{equation}
Let us furthermore set
\begin{equation}
g := \sum_{k\in\mathbb{N}} f_k.
\label{pf.thm.fps.lim.sum-lim.g}
\end{equation}

Now, fix $n\in\mathbb{N}$. Then, all but finitely many
$k\in\mathbb{N}$ satisfy $[x^n] f_k = 0$
(by~\eqref{pf.thm.fps.lim.sum-lim.1}). In other words, there exists a
finite subset $J$ of $\mathbb{N}$ such that
\begin{equation}
\text{all } k\in\mathbb{N}\setminus J \text{ satisfy } [x^n] f_k = 0.
\label{pf.thm.fps.lim.sum-lim.2}
\end{equation}
Consider this subset $J$. The set $J$ is a finite set of nonnegative
integers, and thus has an upper bound. In other words, there exists some
$m\in\mathbb{N}$ such that
\begin{equation}
\text{all } k\in J \text{ satisfy } k\leq m.
\label{pf.thm.fps.lim.sum-lim.3}
\end{equation}
Consider this $m$.

Let $k\in\mathbb{N}$ be such that $k\geq m+1$. If we had $k\in J$,
then we would have $k\leq m$
(by~\eqref{pf.thm.fps.lim.sum-lim.3}), which would contradict
$k\geq m+1 > m$. Thus, we cannot have $k\in J$. Hence,
$k\in\mathbb{N}\setminus J$ (since $k\in\mathbb{N}$ but not
$k\in J$). Therefore,~\eqref{pf.thm.fps.lim.sum-lim.2} yields
$[x^n] f_k = 0$.

Forget that we fixed $k$. We thus have shown that
\begin{equation}
[x^n] f_k = 0 \qquad\text{for each } k\in\mathbb{N} \text{ satisfying }
k\geq m+1.
\label{pf.thm.fps.lim.sum-lim.4}
\end{equation}

Now, let $i$ be an integer such that $i\geq m$. Then,
$i\in\mathbb{N}$ (since $i\geq m\geq 0$) and $i+1\geq m+1$ (since
$i\geq m$). From $g = \sum_{k\in\mathbb{N}} f_k$, we obtain
\begin{align}
[x^n] g &= [x^n]\left(\sum_{k\in\mathbb{N}} f_k\right)
= \sum_{k\in\mathbb{N}} [x^n] f_k \nonumber\\
&= \sum_{k=0}^{i} [x^n] f_k + \sum_{k=i+1}^{\infty}
\underbrace{[x^n] f_k}_{\substack{= 0 \\ \text{(by~\eqref{pf.thm.fps.lim.sum-lim.4}} \\ \text{since } k\geq i+1\geq m+1\text{)}}}
= \sum_{k=0}^{i} [x^n] f_k + \underbrace{\sum_{k=i+1}^{\infty} 0}_{=0} \nonumber\\
&= \sum_{k=0}^{i} [x^n] f_k.
\label{pf.thm.fps.lim.sum-lim.5}
\end{align}
On the other hand, from $g_i = \sum_{k=0}^{i} f_k$, we obtain
\[
[x^n] g_i = [x^n]\left(\sum_{k=0}^{i} f_k\right)
= \sum_{k=0}^{i} [x^n] f_k.
\]
Comparing this with~\eqref{pf.thm.fps.lim.sum-lim.5}, we obtain
$[x^n] g_i = [x^n] g$.

Forget that we fixed $i$. We thus have shown that
\[
\text{all integers } i\geq m \text{ satisfy } [x^n] g_i = [x^n] g.
\]
Hence, there exists some $N\in\mathbb{N}$ such that
\[
\text{all integers } i\geq N \text{ satisfy } [x^n] g_i = [x^n] g
\]
(namely, $N = m$). In other words, the sequence
$\bigl([x^n] g_i\bigr)_{i\in\mathbb{N}}$ stabilizes to $[x^n] g$
(by the definition of ``stabilizes'').

Forget that we fixed $n$. We thus have shown that for each
$n\in\mathbb{N}$,
\[
\text{the sequence } \bigl([x^n] g_i\bigr)_{i\in\mathbb{N}}
\text{ stabilizes to } [x^n] g.
\]
In other words, the sequence $(g_i)_{i\in\mathbb{N}}$ coefficientwise
stabilizes to $g$. In other words, $\lim_{i\to\infty} g_i = g$.
In view of~\eqref{pf.thm.fps.lim.sum-lim.gi}
and~\eqref{pf.thm.fps.lim.sum-lim.g}, we can rewrite this as
\[
\lim_{i\to\infty} \sum_{k=0}^{i} f_k = \sum_{k\in\mathbb{N}} f_k.
\]
Renaming the summation index $k$ as $n$ everywhere in this equality,
we can rewrite it as
\[
\lim_{i\to\infty} \sum_{n=0}^{i} f_n = \sum_{n\in\mathbb{N}} f_n.
\]
This proves Theorem~\ref{thm.fps.lim.sum-lim}.
\end{proof}

\subsection{Helpers for Theorem~\ref{thm.fps.lim.prod-lim}}

\begin{lemma}
\label{lem.fps.lim.coeff-mul-one-plus-higher}
\lean{PowerSeries.coeff_mul_one_plus_higher}
\leanhelper
\leanok
Let $f \in K[[x]]$ satisfy $[x^k] f = \delta_{k,0}$ for all
$k \le n$ (i.e., $f \equiv 1 \pmod{x^{n+1}}$). Then for any
$g \in K[[x]]$ and any $k \le n$,
\[
[x^k](g \cdot f) = [x^k] g.
\]
The proof expands the product coefficient as a sum over the antidiagonal
and observes that only the term with second component $0$ survives,
since $[x^j] f = 0$ for $0 < j \le n$.
\end{lemma}

\begin{lemma}
\label{lem.fps.lim.coeff-prod-extend}
\lean{PowerSeries.coeff_prod_extend}
\leanhelper
\leanok
If $f_m \equiv 1 \pmod{x^{n+1}}$, then for any $k \le n$,
\[
[x^k] \prod_{j=0}^{m} f_j = [x^k] \prod_{j=0}^{m-1} f_j.
\]
That is, extending a partial product by one more term that is
$1 + O(x^{n+1})$ does not change the first $n+1$ coefficients.
\end{lemma}

\begin{lemma}
\label{lem.fps.lim.coeff-prod-range-eq-of-ge}
\lean{PowerSeries.coeff_prod_range_eq_of_ge}
\leanhelper
\leanok
For a multipliable family $(f_j)$, if $f_j \equiv 1 \pmod{x^{n+1}}$
for all $j \ge N$, then for all $i \ge N$ and $k \le n$,
\[
[x^k] \prod_{j=0}^{i} f_j = [x^k] \prod_{j=0}^{N} f_j.
\]
The proof proceeds by induction on $i$, repeatedly applying
Lemma~\ref{lem.fps.lim.coeff-prod-extend}.
\end{lemma}

\begin{lemma}
\label{lem.fps.lim.coeff-partial-prod-stabilizes}
\lean{PowerSeries.coeff_partial_prod_stabilizes}
\leanhelper
\leanok
For a multipliable family $(f_j)$, for each $n \in \mathbb{N}$,
there exists $N$ such that for all $i \ge N$,
\[
[x^n] \prod_{j=0}^{i} f_j = [x^n] \prod_{j\in\mathbb{N}} f_j.
\]
This is the key stabilization lemma that directly implies
Theorem~\ref{thm.fps.lim.prod-lim}.
\end{lemma}

\subsection{Detailed proof of Theorem~\ref{thm.fps.lim.prod-lim}}

\begin{proof}[Proof of Theorem~\ref{thm.fps.lim.prod-lim}]
\proves{thm.fps.lim.prod-lim}
\leanok
The family $(f_n)_{n\in\mathbb{N}}$ is multipliable. In other words,
each coefficient in the product of this family is finitely determined
(by the definition of ``multipliable'').

Let us define
\begin{equation}
g_i := \prod_{k=0}^{i} f_k
\qquad\text{for each } i\in\mathbb{N}.
\label{pf.thm.fps.lim.prod-lim.gi}
\end{equation}
Let us furthermore set
\begin{equation}
g := \prod_{k\in\mathbb{N}} f_k.
\label{pf.thm.fps.lim.prod-lim.g}
\end{equation}

Now, fix $n\in\mathbb{N}$. Then, the $x^n$-coefficient in the product
of the family $(f_k)_{k\in\mathbb{N}}$ is finitely determined. In other
words, there is a finite subset $M$ of $\mathbb{N}$ that determines the
$x^n$-coefficient in the product of $(f_k)_{k\in\mathbb{N}}$. Consider
this subset $M$.

The set $M$ is a finite set of nonnegative integers, and thus has an
upper bound. In other words, there exists some $m\in\mathbb{N}$ such
that
\begin{equation}
\text{all } k\in M \text{ satisfy } k\leq m.
\label{pf.thm.fps.lim.prod-lim.3}
\end{equation}
Consider this $m$.

Now, let $i$ be an integer such that $i\geq m$. Then,
$i\in\mathbb{N}$ (since $i\geq m\geq 0$) and $m\leq i$.
From~\eqref{pf.thm.fps.lim.prod-lim.3}, we see that all $k\in M$
satisfy $k\leq m\leq i$ and therefore $k\in\{0,1,\ldots,i\}$. In other
words, $M\subseteq\{0,1,\ldots,i\}$.

However, the set $M$ determines the $x^n$-coefficient in the product of
$(f_k)_{k\in\mathbb{N}}$. In other words, every finite subset $J$ of
$\mathbb{N}$ satisfying $M\subseteq J\subseteq\mathbb{N}$ satisfies
\[
[x^n]\left(\prod_{k\in J} f_k\right) = [x^n]\left(\prod_{k\in M} f_k\right)
\]
(by the definition of ``determines the $x^n$-coefficient in the product
of $(f_k)_{k\in\mathbb{N}}$''). Applying this to
$J = \{0,1,\ldots,i\}$, we obtain
\begin{equation}
[x^n]\left(\prod_{k\in\{0,1,\ldots,i\}} f_k\right)
= [x^n]\left(\prod_{k\in M} f_k\right)
\label{pf.thm.fps.lim.prod-lim.4}
\end{equation}
(since $\{0,1,\ldots,i\}$ is a finite subset of $\mathbb{N}$ satisfying
$M\subseteq\{0,1,\ldots,i\}\subseteq\mathbb{N}$).

Moreover, the definition of the product of the multipliable family
$(f_k)_{k\in\mathbb{N}}$ shows that
\[
[x^n]\left(\prod_{k\in\mathbb{N}} f_k\right)
= [x^n]\left(\prod_{k\in M} f_k\right)
\]
(since $M$ is a finite subset of $\mathbb{N}$ that determines the
$x^n$-coefficient in the product of $(f_k)_{k\in\mathbb{N}}$).
Comparing this with~\eqref{pf.thm.fps.lim.prod-lim.4}, we find
\[
[x^n]\left(\prod_{k\in\{0,1,\ldots,i\}} f_k\right)
= [x^n]\left(\prod_{k\in\mathbb{N}} f_k\right).
\]
In view of $\prod_{k\in\{0,1,\ldots,i\}} f_k = \prod_{k=0}^{i} f_k = g_i$
(by~\eqref{pf.thm.fps.lim.prod-lim.gi}) and
$\prod_{k\in\mathbb{N}} f_k = g$
(by~\eqref{pf.thm.fps.lim.prod-lim.g}), we can rewrite this as
\[
[x^n] g_i = [x^n] g.
\]

Forget that we fixed $i$. We thus have shown that
\[
\text{all integers } i\geq m \text{ satisfy } [x^n] g_i = [x^n] g.
\]
Hence, the sequence $\bigl([x^n] g_i\bigr)_{i\in\mathbb{N}}$ stabilizes
to $[x^n] g$.

Forget that we fixed $n$. We thus have shown that for each
$n\in\mathbb{N}$, the sequence $\bigl([x^n] g_i\bigr)_{i\in\mathbb{N}}$
stabilizes to $[x^n] g$. In other words, the sequence
$(g_i)_{i\in\mathbb{N}}$ coefficientwise stabilizes to $g$. In other
words, $\lim_{i\to\infty} g_i = g$. In view
of~\eqref{pf.thm.fps.lim.prod-lim.gi}
and~\eqref{pf.thm.fps.lim.prod-lim.g}, we can rewrite this as
\[
\lim_{i\to\infty} \prod_{k=0}^{i} f_k = \prod_{k\in\mathbb{N}} f_k.
\]
Renaming the product index $k$ as $n$ everywhere in this equality, we
can rewrite it as
\[
\lim_{i\to\infty} \prod_{n=0}^{i} f_n = \prod_{n\in\mathbb{N}} f_n.
\]
This proves Theorem~\ref{thm.fps.lim.prod-lim}.
\end{proof}

\subsection{Detailed proofs of Corollary~\ref{cor.fps.lim.fps-as-pol}}

\begin{proof}[Proof of Corollary~\ref{cor.fps.lim.fps-as-pol}]
\proves{cor.fps.lim.fps-as-pol}
\leanok
Let a FPS $a\in K[[x]]$ be written in the form
$a = \sum_{n\in\mathbb{N}} a_n x^n$ with $a_n\in K$.
Then, the family $(a_n x^n)_{n\in\mathbb{N}}$ is summable.
Hence, Theorem~\ref{thm.fps.lim.sum-lim} (applied to $f_n = a_n x^n$)
yields
\[
\lim_{i\to\infty} \sum_{n=0}^{i} a_n x^n
= \sum_{n\in\mathbb{N}} a_n x^n = a.
\]
In other words, $a = \lim_{i\to\infty} \sum_{n=0}^{i} a_n x^n$.
This proves Corollary~\ref{cor.fps.lim.fps-as-pol}.
\end{proof}

\subsection{Detailed proof of Theorem~\ref{thm.fps.lim.sum-lim-conv}}

\begin{proof}[Proof of Theorem~\ref{thm.fps.lim.sum-lim-conv}]
\proves{thm.fps.lim.sum-lim-conv}
\leanok
Let us define
\begin{equation}
g_i := \sum_{k=0}^{i} f_k
\qquad\text{for each } i\in\mathbb{N}.
\label{pf.thm.fps.lim.sum-lim-conv.gi}
\end{equation}
We have assumed that the limit $\lim_{i\to\infty} \sum_{n=0}^{i} f_n$
exists. Let us denote this limit by $g$. Thus,
\begin{equation}
g = \lim_{i\to\infty} \sum_{n=0}^{i} f_n
= \lim_{i\to\infty} g_i.
\label{pf.thm.fps.lim.sum-lim-conv.g=lim}
\end{equation}
In other words, the sequence $(g_i)_{i\in\mathbb{N}}$ coefficientwise
stabilizes to $g$. In other words, for each $n\in\mathbb{N}$,
\begin{equation}
\text{the sequence } \bigl([x^n] g_i\bigr)_{i\in\mathbb{N}}
\text{ stabilizes to } [x^n] g.
\label{pf.thm.fps.lim.sum-lim-conv.gstab}
\end{equation}

Let $n\in\mathbb{N}$. Then, the sequence
$\bigl([x^n] g_i\bigr)_{i\in\mathbb{N}}$ stabilizes to $[x^n] g$
(by~\eqref{pf.thm.fps.lim.sum-lim-conv.gstab}). In other words, there
exists some $N\in\mathbb{N}$ such that
\begin{equation}
\text{all integers } i\geq N \text{ satisfy }
[x^n] g_i = [x^n] g.
\label{pf.thm.fps.lim.sum-lim-conv.1}
\end{equation}
Consider this $N$. Let $M$ be the set $\{0,1,\ldots,N\}$; this is a
finite subset of $\mathbb{N}$.

Let $i\in\mathbb{N}\setminus M$. Thus,
\[
i\in\mathbb{N}\setminus M = \mathbb{N}\setminus\{0,1,\ldots,N\}
= \{N+1, N+2, N+3, \ldots\}.
\]
In other words, $i$ is a nonnegative integer with $i\geq N+1$. Hence,
$i-1\geq N$.
Therefore,~\eqref{pf.thm.fps.lim.sum-lim-conv.1} (applied to $i-1$
instead of $i$) yields $[x^n] g_{i-1} = [x^n] g$.
But~\eqref{pf.thm.fps.lim.sum-lim-conv.1} also yields
$[x^n] g_i = [x^n] g$ (since $i\geq N+1\geq N$). Comparing these two
equalities, we find
\begin{equation}
[x^n] g_{i-1} = [x^n] g_i.
\label{pf.thm.fps.lim.sum-lim-conv.2}
\end{equation}

However,~\eqref{pf.thm.fps.lim.sum-lim-conv.gi} yields
\begin{equation}
g_i = \sum_{k=0}^{i} f_k = f_i + \sum_{k=0}^{i-1} f_k.
\label{pf.thm.fps.lim.sum-lim-conv.3}
\end{equation}
Moreover,~\eqref{pf.thm.fps.lim.sum-lim-conv.gi} (applied to $i-1$
instead of $i$) yields $g_{i-1} = \sum_{k=0}^{i-1} f_k$. In view of
this, we can rewrite~\eqref{pf.thm.fps.lim.sum-lim-conv.3} as
\[
g_i = f_i + g_{i-1}.
\]
Hence,
$[x^n] g_i = [x^n](f_i + g_{i-1}) = [x^n] f_i + \underbrace{[x^n] g_{i-1}}_{\substack{= [x^n] g_i \\ \text{(by~\eqref{pf.thm.fps.lim.sum-lim-conv.2})}}}
= [x^n] f_i + [x^n] g_i$.
Subtracting $[x^n] g_i$ from both sides of this equality, we find
$0 = [x^n] f_i$. In other words, $[x^n] f_i = 0$.

Forget that we fixed $i$. We thus have shown that all
$i\in\mathbb{N}\setminus M$ satisfy $[x^n] f_i = 0$. Hence, all but
finitely many $i\in\mathbb{N}$ satisfy $[x^n] f_i = 0$ (since all but
finitely many $i\in\mathbb{N}$ belong to $\mathbb{N}\setminus M$
because the set $M$ is finite). Renaming the index $i$ as $k$ in this
statement, we obtain the following: All but finitely many
$k\in\mathbb{N}$ satisfy $[x^n] f_k = 0$.

Forget that we fixed $n$. We thus have shown that for each
$n\in\mathbb{N}$,
\[
\text{all but finitely many } k\in\mathbb{N} \text{ satisfy }
[x^n] f_k = 0.
\]
In other words, the family $(f_n)_{n\in\mathbb{N}}$ is summable (by the
definition of ``summable'').

Hence, Theorem~\ref{thm.fps.lim.sum-lim} yields
$\lim_{i\to\infty} \sum_{n=0}^{i} f_n = \sum_{n\in\mathbb{N}} f_n$.
In other words,
$\sum_{n\in\mathbb{N}} f_n = \lim_{i\to\infty} \sum_{n=0}^{i} f_n$.
This completes the proof of Theorem~\ref{thm.fps.lim.sum-lim-conv}.
\end{proof}

\subsection{Helpers for Theorem~\ref{thm.fps.lim.prod-lim-conv}}

\begin{lemma}
\label{lem.fps.lim.constantCoeff-prod-eq-one}
\lean{PowerSeries.constantCoeff_prod_eq_one}
\leanhelper
\leanok
If $f_i$ has constant term $1$ for every $i$ in a finite set $S$,
then $\prod_{i \in S} f_i$ also has constant term $1$.
\end{lemma}

\subsection{Detailed proof of Theorem~\ref{thm.fps.lim.prod-lim-conv}}

\begin{proof}[Proof of Theorem~\ref{thm.fps.lim.prod-lim-conv}]
\proves{thm.fps.lim.prod-lim-conv}
\leanok
Let us define
\begin{equation}
g_i := \prod_{k=0}^{i} f_k
\qquad\text{for each } i\in\mathbb{N}.
\label{pf.thm.fps.lim.prod-lim-conv.gi}
\end{equation}

We have assumed that the limit $\lim_{i\to\infty} \prod_{n=0}^{i} f_n$
exists. Let us denote this limit by $g$. Thus,
\begin{equation}
g = \lim_{i\to\infty} \prod_{n=0}^{i} f_n = \lim_{i\to\infty} g_i.
\label{pf.thm.fps.lim.prod-lim-conv.g=lim}
\end{equation}

Let $n\in\mathbb{N}$ be arbitrary.
Lemma~\ref{lem.fps.lim.xn-equiv} (applied to $g_i$ and $g$ instead of
$f_i$ and $f$) shows that there exists some integer $N\in\mathbb{N}$
such that
\begin{equation}
\text{all integers } i\geq N \text{ satisfy }
g_i \overset{x^n}{\equiv} g
\label{pf.thm.fps.lim.prod-lim-conv.gi-equiv-g}
\end{equation}
(since $g = \lim_{i\to\infty} g_i$). Consider this $N$.

Let $M = \{0,1,\ldots,N\}$. Thus, $M$ is a finite subset of
$\mathbb{N}$. We shall show that this subset $M$ determines the
$x^n$-coefficient in the product of $(f_k)_{k\in\mathbb{N}}$.

For this purpose, we first observe that $M = \{0,1,\ldots,N\}$, so that
\begin{equation}
\prod_{k\in M} f_k = \prod_{k\in\{0,1,\ldots,N\}} f_k
= \prod_{k=0}^{N} f_k = g_N
\label{pf.thm.fps.lim.prod-lim-conv.=gN}
\end{equation}
(since $g_N$ was defined to be $\prod_{k=0}^{N} f_k$).
Applying~\eqref{pf.thm.fps.lim.prod-lim-conv.gi-equiv-g} to $i = N$,
we obtain $g_N \overset{x^n}{\equiv} g$ (since $N\geq N$). Therefore,
$g \overset{x^n}{\equiv} g_N$ (since $\overset{x^n}{\equiv}$ is an
equivalence relation and thus symmetric).

Next, we shall show two claims:

\textbf{Claim 1:} \textit{For each integer $j > N$, we have
$g \overset{x^n}{\equiv} g f_j$.}

\begin{proof}[Proof of Claim 1]
\leanok
Let $j > N$ be an integer. Thus, $j\geq N+1$, so that $j-1\geq N$.
Hence,~\eqref{pf.thm.fps.lim.prod-lim-conv.gi-equiv-g} (applied to
$i = j-1$) yields $g_{j-1} \overset{x^n}{\equiv} g$.
However,~\eqref{pf.thm.fps.lim.prod-lim-conv.gi-equiv-g} (applied to
$i = j$) yields $g_j \overset{x^n}{\equiv} g$ (since
$j\geq N+1\geq N$). Thus, $g \overset{x^n}{\equiv} g_j$ (by symmetry).

Also, we obviously have $f_j \overset{x^n}{\equiv} f_j$ (by
reflexivity).

However, the definition of $g_{j-1}$ yields
$g_{j-1} = \prod_{k=0}^{j-1} f_k$. Furthermore, the definition of
$g_j$ yields
\[
g_j = \prod_{k=0}^{j} f_k
= \left(\prod_{k=0}^{j-1} f_k\right) f_j = g_{j-1} f_j.
\]
But recall that $g_{j-1} \overset{x^n}{\equiv} g$ and
$f_j \overset{x^n}{\equiv} f_j$. Hence,
$g_{j-1} f_j \overset{x^n}{\equiv} g f_j$ (by the multiplicativity of
$x^n$-equivalence). In view of $g_j = g_{j-1} f_j$, we can rewrite
this as $g_j \overset{x^n}{\equiv} g f_j$.

Now, recall that $g \overset{x^n}{\equiv} g_j$. Thus,
$g \overset{x^n}{\equiv} g_j \overset{x^n}{\equiv} g f_j$, so that
$g \overset{x^n}{\equiv} g f_j$ (by transitivity). This proves Claim 1.
\end{proof}

\textbf{Claim 2:} \textit{If $U$ is any finite subset of $\mathbb{N}$
satisfying $M\subseteq U$, then
$g \overset{x^n}{\equiv} \prod_{k\in U} f_k$.}

\begin{proof}[Proof of Claim 2]
\leanok
We induct on the nonnegative integer $|U\setminus M|$:

\textit{Base case:} Let us check that Claim 2 holds when
$|U\setminus M| = 0$.

Indeed, let $U$ be any finite subset of $\mathbb{N}$ satisfying
$M\subseteq U$ and $|U\setminus M| = 0$. Then,
$U\setminus M = \varnothing$ (since $|U\setminus M| = 0$), so that
$U\subseteq M$. Combined with $M\subseteq U$, this yields $U = M$.
Thus, $\prod_{k\in U} f_k = \prod_{k\in M} f_k = g_N$
(by~\eqref{pf.thm.fps.lim.prod-lim-conv.=gN}). But we have already
proved that $g \overset{x^n}{\equiv} g_N$. In other words,
$g \overset{x^n}{\equiv} \prod_{k\in U} f_k$ (since
$\prod_{k\in U} f_k = g_N$). This completes the base case.

\textit{Induction step:} Let $s\in\mathbb{N}$. Assume (as the induction
hypothesis) that Claim 2 holds when $|U\setminus M| = s$. We must now
prove that Claim 2 holds when $|U\setminus M| = s+1$.

Indeed, let $U$ be any finite subset of $\mathbb{N}$ satisfying
$M\subseteq U$ and $|U\setminus M| = s+1$. Then, the set
$U\setminus M$ is nonempty (since $|U\setminus M| = s+1 > 0$). Hence,
there exists some $u\in U\setminus M$. Consider this $u$. By the
induction hypothesis (applied to $U\setminus\{u\}$ instead of $U$), we
obtain $g \overset{x^n}{\equiv} \prod_{k\in U\setminus\{u\}} f_k$.

However, $u\in U\setminus M$, so $u\notin M = \{0,1,\ldots,N\}$, which
gives $u > N$. Hence, Claim 1 (applied to $j = u$) yields
$g \overset{x^n}{\equiv} g f_u$.

We have $g \overset{x^n}{\equiv} \prod_{k\in U\setminus\{u\}} f_k$ and
$f_u \overset{x^n}{\equiv} f_u$. Hence,
$g f_u \overset{x^n}{\equiv} \left(\prod_{k\in U\setminus\{u\}} f_k\right) f_u$
(by the multiplicativity of $x^n$-equivalence). In view of
$\prod_{k\in U} f_k = \left(\prod_{k\in U\setminus\{u\}} f_k\right) f_u$,
we can rewrite this as
$g f_u \overset{x^n}{\equiv} \prod_{k\in U} f_k$.

Now, recall that $g \overset{x^n}{\equiv} g f_u$. Thus,
$g \overset{x^n}{\equiv} g f_u \overset{x^n}{\equiv} \prod_{k\in U} f_k$,
so $g \overset{x^n}{\equiv} \prod_{k\in U} f_k$ (by transitivity).
This completes the induction step. Thus, Claim 2 is proved by induction.
\end{proof}

Now, let $J$ be a finite subset of $\mathbb{N}$ satisfying
$M\subseteq J\subseteq\mathbb{N}$. Then,
$g \overset{x^n}{\equiv} \prod_{k\in J} f_k$ (by Claim 2, applied to
$U = J$). In other words,
\[
\text{each } m\in\{0,1,\ldots,n\} \text{ satisfies }
[x^m] g = [x^m]\left(\prod_{k\in J} f_k\right)
\]
(by the definition of $x^n$-equivalence). Applying this to $m = n$, we
find
\[
[x^n] g = [x^n]\left(\prod_{k\in J} f_k\right).
\]
The same argument can be applied to $M$ instead of $J$ (since
$M\subseteq M\subseteq\mathbb{N}$), and thus yields
\[
[x^n] g = [x^n]\left(\prod_{k\in M} f_k\right).
\]
Comparing these two equalities, we find
\[
[x^n]\left(\prod_{k\in J} f_k\right)
= [x^n]\left(\prod_{k\in M} f_k\right).
\]

Forget that we fixed $J$. We have now shown that every finite subset $J$
of $\mathbb{N}$ satisfying $M\subseteq J\subseteq\mathbb{N}$ satisfies
\[
[x^n]\left(\prod_{k\in J} f_k\right)
= [x^n]\left(\prod_{k\in M} f_k\right).
\]
In other words, the set $M$ determines the $x^n$-coefficient in the
product of $(f_k)_{k\in\mathbb{N}}$. Hence, the $x^n$-coefficient in
the product of $(f_k)_{k\in\mathbb{N}}$ is finitely determined (since
$M$ is a finite subset of $\mathbb{N}$).

Forget that we fixed $n$. We thus have shown that for each
$n\in\mathbb{N}$, the $x^n$-coefficient in the product of
$(f_k)_{k\in\mathbb{N}}$ is finitely determined. In other words, each
coefficient in the product of $(f_k)_{k\in\mathbb{N}}$ is finitely
determined. In other words, the family $(f_k)_{k\in\mathbb{N}}$ is
multipliable. In other words, the family $(f_n)_{n\in\mathbb{N}}$ is
multipliable.

Hence, Theorem~\ref{thm.fps.lim.prod-lim} yields
$\lim_{i\to\infty} \prod_{n=0}^{i} f_n = \prod_{n\in\mathbb{N}} f_n$.
In other words,
$\prod_{n\in\mathbb{N}} f_n = \lim_{i\to\infty} \prod_{n=0}^{i} f_n$.
This completes the proof of Theorem~\ref{thm.fps.lim.prod-lim-conv}.
\end{proof}

\subsection{Helper lemmas from the Lean formalization}

\begin{lemma}
\label{lem.fps.lim.stab-unique-detail}
\lean{Seq.stabilizesTo_unique'}
\leanhelper
\leanok
If a sequence $(a_i)_{i\in\mathbb{N}}$ in $K$ stabilizes to both
$\ell_1$ and $\ell_2$, then $\ell_1 = \ell_2$.
\end{lemma}

\begin{proof}
\leanok
Let $N_1$ and $N_2$ be the stabilization indices for $\ell_1$ and
$\ell_2$ respectively. Set $N := \max(N_1, N_2)$. Then
$a_N = \ell_1$ and $a_N = \ell_2$, whence $\ell_1 = \ell_2$.
\end{proof}

\begin{lemma}
\label{lem.fps.lim.partial-sum-detail}
\lean{PowerSeries.coeffStabilizesTo_partial_sum'}
\leanhelper
\leanok
If the family $(f_n)_{n\in\mathbb{N}}$ is summable, then the sequence of
partial sums $\left(\sum_{j=0}^{i} f_j\right)_{i\in\mathbb{N}}$
coefficientwise stabilizes to $\sum_{n\in\mathbb{N}} f_n$.
\end{lemma}

\begin{proof}
\leanok
This is a direct restatement of Theorem~\ref{thm.fps.lim.sum-lim}.
\end{proof}

\begin{lemma}
\label{lem.fps.lim.partial-prod-detail}
\lean{PowerSeries.coeffStabilizesTo_partial_prod'}
\leanhelper
\leanok
If the family $(f_n)_{n\in\mathbb{N}}$ is multipliable, then the
sequence of partial products
$\left(\prod_{j=0}^{i} f_j\right)_{i\in\mathbb{N}}$ coefficientwise
stabilizes to $\prod_{n\in\mathbb{N}} f_n$.
\end{lemma}

\begin{proof}
\leanok
This is a direct restatement of Theorem~\ref{thm.fps.lim.prod-lim}.
\end{proof}

\begin{lemma}
\label{lem.fps.lim.summable-of-partial-sum-detail}
\lean{PowerSeries.isSummable_of_coeffStabilizesTo_partial_sum'}
\leanhelper
\leanok
If the sequence of partial sums
$\left(\sum_{j=0}^{i} f_j\right)_{i\in\mathbb{N}}$ coefficientwise
stabilizes to some limit, then the family $(f_n)_{n\in\mathbb{N}}$ is
summable.
\end{lemma}

\begin{proof}
\leanok
This is a direct restatement of the summability part of
Theorem~\ref{thm.fps.lim.sum-lim-conv}.
\end{proof}

\begin{lemma}
\label{lem.fps.lim.tsum-eq-detail}
\lean{PowerSeries.tsum'_eq_of_coeffStabilizesTo_partial_sum'}
\leanhelper
\leanok
If the sequence of partial sums
$\left(\sum_{j=0}^{i} f_j\right)_{i\in\mathbb{N}}$ coefficientwise
stabilizes to some limit and the family $(f_n)_{n\in\mathbb{N}}$ is
summable, then the infinite sum equals the limit.
\end{lemma}

\begin{proof}
\leanok
By Theorem~\ref{thm.fps.lim.sum-lim}, the partial sums converge to the
infinite sum. By uniqueness of limits
(Lemma~\ref{lem.fps.lim.coeff-stab-unique}), the limit equals the
infinite sum.
\end{proof}

\begin{lemma}
\label{lem.fps.lim.multipliable-of-partial-prod-detail}
\lean{PowerSeries.isMultipliable_of_coeffStabilizesTo_partial_prod'}
\leanhelper
\leanok
If the sequence of partial products
$\left(\prod_{j=0}^{i} f_j\right)_{i\in\mathbb{N}}$ coefficientwise
stabilizes to some limit, then the
family $(f_n)_{n\in\mathbb{N}}$ is multipliable.
\end{lemma}

\begin{proof}
\leanok
This is a direct restatement of the multipliability part of
Theorem~\ref{thm.fps.lim.prod-lim-conv}.
\end{proof}

\begin{lemma}
\label{lem.fps.lim.tprod-eq-detail}
\lean{PowerSeries.tprod'_eq_of_coeffStabilizesTo_partial_prod'}
\leanhelper
\leanok
If the sequence of partial products
$\left(\prod_{j=0}^{i} f_j\right)_{i\in\mathbb{N}}$ coefficientwise
stabilizes to some limit and the family $(f_n)_{n\in\mathbb{N}}$ is
multipliable, then the infinite product equals the limit.
\end{lemma}

\begin{proof}
\leanok
By Theorem~\ref{thm.fps.lim.prod-lim}, the partial products converge to
the infinite product. By uniqueness of limits
(Lemma~\ref{lem.fps.lim.coeff-stab-unique}), the limit equals the
infinite product.
\end{proof}

\begin{lemma}
\label{lem.fps.lim.xnEquiv-pow}
\lean{PowerSeries.xnEquiv_pow}
\leanhelper
\leanok
If $f \overset{x^n}{\equiv} g$, then $f^d \overset{x^n}{\equiv} g^d$
for every $d \in \mathbb{N}$. The proof proceeds by induction on $d$,
using the multiplicativity of $x^n$-equivalence at each step.
\end{lemma}
